\documentclass[11pt]{article}
\usepackage[margin=1in]{geometry}
\usepackage[T1]{fontenc}
\usepackage{lmodern}
\usepackage{enumitem}

\title{ICSPSET Problem-wise Change Explainer}
\author{}
\date{}

\begin{document}
\maketitle

\section*{Source Mapping}
\begin{itemize}[leftmargin=1.2em]
    \item \textbf{Problem statements}: \texttt{Instructions \& Problem Set}
    \item \textbf{Code base}: \texttt{ics\_images.py}
\end{itemize}

\section*{Problem 2-1: Socialize with other pixels}
\subsection*{(a) Upgrade 3\(\times\)3 convolution to 5\(\times\)5}
\textbf{What changes are made:}
\begin{itemize}[leftmargin=1.2em]
    \item Extend convolution neighborhood logic from radius \(1\) to radius \(2\).
    \item Replace fixed 3\(\times\)3 channel accumulation in \texttt{multiply(...)} with 5\(\times\)5 accumulation.
\end{itemize}
\textbf{Why:} The objective explicitly asks for a 5\(\times\)5 kernel implementation using the given codebase.\\
\textbf{How this meets objective:} Each output pixel is computed from a 25-pixel neighborhood, so filtering behavior matches the required 5\(\times\)5 convolution.

\subsection*{(b) Generalize to any odd \(n>1\)}
\textbf{What changes are made:}
\begin{itemize}[leftmargin=1.2em]
    \item Parameterize kernel size as \(n\), with radius \(r=\lfloor n/2 \rfloor\).
    \item Add input validation: reject even \(n\), and reject \(n \le 1\), with a clear error message.
    \item Use loop-based accumulation over the \(n \times n\) kernel instead of hardcoded indices.
\end{itemize}
\textbf{Why:} The task requires scalability beyond 5\(\times\)5 while enforcing odd-sized kernels only.\\
\textbf{How this meets objective:} Dynamic indexing computes valid odd-kernel convolutions for all \(n>1\), and invalid sizes fail fast with explicit errors.

\section*{Problem 2-2: Take a left to go right}
\subsection*{(a) Rotate clockwise by \(90^\circ\)}
\textbf{What changes are made:}
\begin{itemize}[leftmargin=1.2em]
    \item Add a function that maps each input pixel \((i,j)\) to clockwise output coordinates.
    \item Swap output dimensions from \((H,W)\) to \((W,H)\).
\end{itemize}
\textbf{Why:} A \(90^\circ\) rotation is a geometric transform, not a color transform; it must reindex pixels.\\
\textbf{How this meets objective:} The output image orientation is rotated right exactly once with preserved pixel values.

\subsection*{(b) Rotate counter-clockwise by \(90^\circ\)}
\textbf{What changes are made:}
\begin{itemize}[leftmargin=1.2em]
    \item Add the inverse coordinate mapping for a left \(90^\circ\) rotation.
    \item Keep the same output dimension swap logic.
\end{itemize}
\textbf{Why:} Counter-clockwise rotation is separately required and uses a different index transform.\\
\textbf{How this meets objective:} The output image orientation is rotated left exactly once with no pixel loss.

\section*{Problem 2-3: Reshaping reality (kind of)}
\subsection*{(a) Crop using point \(P=(x,y)\), width \(w\), height \(h\)}
\textbf{What changes are made:}
\begin{itemize}[leftmargin=1.2em]
    \item Add crop logic that extracts rows \(y \rightarrow y+h-1\) and columns \(x \rightarrow x+w-1\).
    \item Return a new image region with dimensions \(w \times h\).
\end{itemize}
\textbf{Why:} Cropping isolates a rectangular region defined by top-left anchor and size.\\
\textbf{How this meets objective:} The generated image is exactly the requested spatial subset.

\subsection*{(b) Stretch/squeeze to \texttt{displayW}, \texttt{displayH}}
\textbf{What changes are made:}
\begin{itemize}[leftmargin=1.2em]
    \item Add resize logic that remaps cropped-region coordinates into requested display dimensions.
    \item Use deterministic coordinate scaling so every destination pixel samples from valid source coordinates.
\end{itemize}
\textbf{Why:} The task requires changing display dimensions after cropping.\\
\textbf{How this meets objective:} The cropped image is transformed from \((w,h)\) to \((displayW,displayH)\), implementing stretch (upscale) or squeeze (downscale).

\subsection*{(c) Safety checks / edge cases}
\textbf{Checks to implement and document:}
\begin{itemize}[leftmargin=1.2em]
    \item Kernel checks: \(n\) is odd and \(n>1\), and kernel dimensions match \(n\times n\).
    \item Image bounds: reject crop start outside image; reject non-positive \(w,h\); prevent out-of-range crop extents.
    \item Resize checks: reject non-positive \texttt{displayW}, \texttt{displayH}.
    \item Type checks: ensure integer dimensions and indices where required.
    \item Border handling consistency during convolution (copy pixel, pad, or other chosen policy) should be explicit and consistent.
\end{itemize}
\textbf{Why:} These validations prevent invalid indexing, runtime errors, and ambiguous outputs.\\
\textbf{How this meets objective:} Robust input handling ensures output generation remains correct, predictable, and aligned with assignment constraints.

\end{document}
